\chapter*{คำนำ}
\addcontentsline{toc}{chapter}{คำนำ}

ฟิสิกส์เป็นศาสตร์พื้นฐานอันเป็นรากฐานของวิทยาศาสตร์และเทคโนโลยีสมัยใหม่ การทำความเข้าใจปรากฏการณ์ธรรมชาติอย่างลึกซึ้ง จำเป็นต้องอาศัยการบูรณาการอย่างกลมกลืนระหว่างแบบจำลองทางคณิตศาสตร์ที่รัดกุม การทดลองตรวจวัดที่แม่นยำ และระเบียบวิธีวิจัยเชิงคำนวณที่ทันสมัย

เอกสารคำสอนและคู่มือวิจัยเรื่อง \textbf{เอกสารคำสอนและคู่มือวิจัยฟิสิกส์ระดับมหาวิทยาลัย} เล่มนี้ ได้รับการพัฒนาขึ้นโดยสังเคราะห์มาตรฐานการจัดทำตำราวิชาการระดับโลกจากตำราชั้นนำสองเล่ม ได้แก่ กรอบการแก้ปัญหาทางฟิสิกส์แบบสี่ขั้นตอน P-S-C-T (Picture, Strategy, Calculation, Test) ของ Tipler \& Mosca และระเบียบวิธีวิเคราะห์โจทย์ปัญหาแบบ C-C-A-F (Conceptualize, Categorize, Analyze, Finalize) ของ Serway \& Jewett เพื่อเสริมสร้างกระบวนการคิดเชิงวิเคราะห์และการแก้ปัญหาอย่างเป็นระบบให้แก่นักศึกษาสาขาวิชาฟิสิกส์และวิทยาศาสตร์กายภาพ

โครงสร้างเนื้อหาครอบคลุม 8 บทเรียนสำคัญ ตั้งแต่การวิเคราะห์ความไม่แน่นอนตามมาตรฐานสากล ISO/IEC Guide 98-3 (GUM), พลศาสตร์ขั้นสูงและสมการการเคลื่อนที่แบบไม่เชิงเส้น, ทฤษฎีศักย์และสมดุลเชิงเสถียรภาพ, ทฤษฎีสนามแม่เหล็กไฟฟ้าและสมการแมกซ์เวลล์ในรูปอนุพันธ์ย่อย, ทัศนศาสตร์เชิงคลื่นและการแทรกสอดของแสงเลเซอร์, กลศาสตร์ควอนตัมและแบบจำลองอะตอม, ฟิสิกส์ของแข็งและวัสดุแม่เหล็กขั้นสูง ซึ่งบูรณาการผลงานวิจัยจริงของผู้เขียนด้านการจำลองโครงสร้างอิเล็กตรอน $\text{LSDA}+U$ ในสารประกอบ $\text{MnO}$ และ $\text{NiO}$, ไปจนถึงระเบียบวิธีวิจัยและฟิสิกส์เชิงคำนวณด้วย Python และ Physics-Informed Neural Networks (PINNs)

ผู้เขียนหวังเป็นอย่างยิ่งว่า เอกสารคำสอนเล่มนี้จะไม่เพียงแต่นำเสนอองค์ความรู้ทางทฤษฎีเท่านั้น แต่ยังทำหน้าที่เป็นสะพานเชื่อมโยงนักศึกษาจากการเรียนรู้ในห้องเรียนสู่โลกแห่งการวิจัยฟิสิกส์ระดับสากลได้อย่างสมภาคภูมิ

\vspace{1.5cm}
\begin{flushright}
\textbf{ผศ.ดร.ชีวะ ทัศนา}\\
สาขาวิชาฟิสิกส์ คณะวิทยาศาสตร์และเทคโนโลยี\\
มหาวิทยาลัยราชภัฏรำไพพรรณี
\end{flushright}
